% Generated by analysis/latex_export.py -- do not edit by hand.
% % Generated by analysis/latex_export.py -- do not edit by hand.
% % Generated by analysis/latex_export.py -- do not edit by hand.
% % Generated by analysis/latex_export.py -- do not edit by hand.
% \input{macros} in the preamble; every macro is prefixed with \scale.
% Each line ends with a comment saying what the value means, measures or represents
% (descriptions live in analysis/macro_docs.py). Percentages are relative differences
% unless stated; 'design level' means one value per mix design, repeats averaged.
\providecommand{\xspace}{}
\newcommand{\scaleNEvaluations}{7{,}944\xspace}                          % number of MoST pass/fail SLA evaluations analysed (all runs, after removing exact copies)
\newcommand{\scaleNMixRuns}{301\xspace}                                  % number of mix runs (MoST searches on request mixes)
\newcommand{\scaleNDesigns}{114\xspace}                                  % number of distinct mix designs (model x test x composition); unit of the frequentist tests
\newcommand{\scaleNRepeatRuns}{270\xspace}                               % number of mix runs that belong to designs measured more than once
\newcommand{\scaleNRepeatDesigns}{83\xspace}                             % number of mix designs measured more than once (stability check)
\newcommand{\scaleNDuplicates}{114\xspace}                               % number of byte-identical copies of a run found in the archive and counted once
\newcommand{\scaleNPureRuns}{360\xspace}                                 % number of pure-class runs (one request class per MoST search)
\newcommand{\scaleQHatFreq}{\ensuremath{-}0.91\xspace}                   % power-mean exponent q of the combination law, profile-likelihood estimate (q=-1 harmonic, 0 geometric, +1 arithmetic)
\newcommand{\scaleQFreqLo}{\ensuremath{-}0.98\xspace}                    % lower bound of the 95 % profile-likelihood confidence interval of q
\newcommand{\scaleQFreqHi}{\ensuremath{-}0.84\xspace}                    % upper bound of the 95 % profile-likelihood confidence interval of q
\newcommand{\scaleHarmonicWinShare}{97\,\%\xspace}                       % share of mix runs for which the harmonic prediction is closer to the measurement than the arithmetic one
\newcommand{\scaleCalibSlope}{0.96\xspace}                               % slope of the regression log(measured mix capacity) ~ log(harmonic prediction); 1 = correct scaling
\newcommand{\scaleCalibSlopeLo}{0.90\xspace}                             % lower bound of the 95 % confidence interval of the calibration slope (cluster-robust by model)
\newcommand{\scaleCalibSlopeHi}{1.02\xspace}                             % upper bound of the 95 % confidence interval of the calibration slope
\newcommand{\scaleCalibRsq}{0.97\xspace}                                 % R^2 of the calibration regression on the log scale
\newcommand{\scaleBiasFreq}{\ensuremath{-}3.0\,\%\xspace}                % mean departure of measured mix capacity from the harmonic prediction, frequentist, designs as units
\newcommand{\scaleBiasFreqLo}{\ensuremath{-}6.9\,\%\xspace}              % lower bound of the 95 % confidence interval of BiasFreq
\newcommand{\scaleBiasFreqHi}{+1.0\,\%\xspace}                           % upper bound of the 95 % confidence interval of BiasFreq
\newcommand{\scaleBiasTostP}{0.001\xspace}                               % p-value of the TOST equivalence test of the bias against +-10 % (small = practically zero)
\newcommand{\scaleArithMeanErr}{\ensuremath{-}53\,\%\xspace}             % mean error of the arithmetic-mean law (measured vs predicted mix capacity, design level); negative = overestimates
\newcommand{\scaleGeomMeanErr}{\ensuremath{-}39\,\%\xspace}              % mean error of the geometric-mean law (design level); negative = overestimates
\newcommand{\scaleBoundsIn}{294\xspace}                                  % number of mix runs whose capacity lies between their slowest and fastest class
\newcommand{\scaleBoundsN}{301\xspace}                                   % number of mix runs checked for boundedness
\newcommand{\scaleBoundsShare}{97.7\,\%\xspace}                          % share of mix runs within the capacity range of their classes
\newcommand{\scaleAboveArith}{0\xspace}                                  % number of mix runs whose capacity exceeds the arithmetic mean of their classes
\newcommand{\scaleQHatBayes}{\ensuremath{-}0.96\xspace}                  % power-mean exponent q, Bayesian posterior median (free-q model)
\newcommand{\scaleQBayesLo}{\ensuremath{-}1.02\xspace}                   % lower bound of the 95 % credible interval of q
\newcommand{\scaleQBayesHi}{\ensuremath{-}0.88\xspace}                   % upper bound of the 95 % credible interval of q
\newcommand{\scaleBiasBayes}{\ensuremath{-}5.5\,\%\xspace}               % pooled bias delta_0 of mixes w.r.t. the harmonic law, Bayesian posterior median
\newcommand{\scaleBiasBayesLo}{\ensuremath{-}9.9\,\%\xspace}             % lower bound of the 90 % credible interval of delta_0
\newcommand{\scaleBiasBayesHi}{\ensuremath{-}1.0\,\%\xspace}             % upper bound of the 90 % credible interval of delta_0
\newcommand{\scaleBiasBayesPRope}{92\,\%\xspace}                         % posterior probability that |delta_0| < 10 % (region of practical equivalence)
\newcommand{\scalePredFree}{+1.7\xspace}                                 % out-of-sample log predictive density of the free-q law minus that of the harmonic law (higher = better)
\newcommand{\scalePredFreeSE}{3.3\xspace}                                % standard error of PredFree
\newcommand{\scalePredGeom}{\ensuremath{-}264.5\xspace}                  % out-of-sample log predictive density of the geometric law minus that of the harmonic law
\newcommand{\scalePredGeomSE}{24.8\xspace}                               % standard error of PredGeom
\newcommand{\scalePredArith}{\ensuremath{-}441.4\xspace}                 % out-of-sample log predictive density of the arithmetic law minus that of the harmonic law
\newcommand{\scalePredArithSE}{38.1\xspace}                              % standard error of PredArith
\newcommand{\scaleSdMixMIT}{6\,\%\xspace}                                % posterior SD of the design effect (departure of a mix design from the harmonic law), MIT (2-minute instantaneous test), as %
\newcommand{\scaleSdRunMIT}{0.3\,\%\xspace}                              % posterior SD between repeated runs of the same mix design, MIT (2-minute instantaneous test), as %
\newcommand{\scaleSdRepMIT}{0.3\,\%\xspace}                              % posterior SD between repeated runs of the same pure class, MIT (2-minute instantaneous test), as %
\newcommand{\scaleBandMIT}{\ensuremath{\pm}15\,\%\xspace}                % half-width of the 90 % planning band around the harmonic prediction, MIT (2-minute instantaneous test)
\newcommand{\scaleSdMixMST}{21\,\%\xspace}                               % posterior SD of the design effect (departure of a mix design from the harmonic law), MST (30-minute sustainable test), as %
\newcommand{\scaleSdRunMST}{1.8\,\%\xspace}                              % posterior SD between repeated runs of the same mix design, MST (30-minute sustainable test), as %
\newcommand{\scaleSdRepMST}{1.7\,\%\xspace}                              % posterior SD between repeated runs of the same pure class, MST (30-minute sustainable test), as %
\newcommand{\scaleBandMST}{\ensuremath{\pm}39\,\%\xspace}                % half-width of the 90 % planning band around the harmonic prediction, MST (30-minute sustainable test)
\newcommand{\scaleBiasDeepSeekMIT}{\ensuremath{-}8.8\,\%\xspace}         % Bayesian bias delta of DeepSeek-R1-Distill-Qwen-7B mixes w.r.t. the harmonic law, MIT (2-minute instantaneous test)
\newcommand{\scaleBiasDeepSeekMST}{\ensuremath{-}4.1\,\%\xspace}         % Bayesian bias delta of DeepSeek-R1-Distill-Qwen-7B mixes w.r.t. the harmonic law, MST (30-minute sustainable test)
\newcommand{\scaleBiasGemmaMIT}{\ensuremath{-}0.3\,\%\xspace}            % Bayesian bias delta of Gemma-7B mixes w.r.t. the harmonic law, MIT (2-minute instantaneous test)
\newcommand{\scaleBiasGemmaMST}{\ensuremath{-}4.6\,\%\xspace}            % Bayesian bias delta of Gemma-7B mixes w.r.t. the harmonic law, MST (30-minute sustainable test)
\newcommand{\scaleBiasLlamaMIT}{\ensuremath{-}7.0\,\%\xspace}            % Bayesian bias delta of Llama-3.1-8B-Instruct mixes w.r.t. the harmonic law, MIT (2-minute instantaneous test)
\newcommand{\scaleBiasLlamaMST}{\ensuremath{-}7.8\,\%\xspace}            % Bayesian bias delta of Llama-3.1-8B-Instruct mixes w.r.t. the harmonic law, MST (30-minute sustainable test)
\newcommand{\scaleBayesDraws}{600\xspace}                                % NUTS draws per chain (4 chains, same number of tuning steps)
\newcommand{\scaleMaxRhat}{1.02\xspace}                                  % largest split-R-hat over the global parameters of the reported Bayesian models (<1.05 = converged)
\newcommand{\scaleMaeMIT}{10\,\%\xspace}                                 % mean absolute error of the harmonic prediction per mix design, MIT (2-minute instantaneous test)
\newcommand{\scaleBiasDesignMIT}{\ensuremath{-}2.9\,\%\xspace}           % mean departure from the harmonic law per mix design, MIT (2-minute instantaneous test)
\newcommand{\scaleWithinTenMIT}{61\,\%\xspace}                           % share of mix designs predicted within +-10 %, MIT (2-minute instantaneous test)
\newcommand{\scaleWithinTenMITLight}{75\,\%\xspace}                      % share of mix designs without the heaviest class predicted within +-10 %, MIT (2-minute instantaneous test)
\newcommand{\scaleWithinTenMITHeavy}{42\,\%\xspace}                      % share of mix designs with the heaviest class predicted within +-10 %, MIT (2-minute instantaneous test)
\newcommand{\scaleMaeMST}{23\,\%\xspace}                                 % mean absolute error of the harmonic prediction per mix design, MST (30-minute sustainable test)
\newcommand{\scaleBiasDesignMST}{\ensuremath{-}3.1\,\%\xspace}           % mean departure from the harmonic law per mix design, MST (30-minute sustainable test)
\newcommand{\scaleWithinTenMST}{48\,\%\xspace}                           % share of mix designs predicted within +-10 %, MST (30-minute sustainable test)
\newcommand{\scaleWithinTenMSTLight}{58\,\%\xspace}                      % share of mix designs without the heaviest class predicted within +-10 %, MST (30-minute sustainable test)
\newcommand{\scaleWithinTenMSTHeavy}{25\,\%\xspace}                      % share of mix designs with the heaviest class predicted within +-10 %, MST (30-minute sustainable test)
\newcommand{\scaleMitMstMeanP}{0.507\xspace}                             % p-value (Mann-Whitney) for a difference in mean bias between MIT and MST designs
\newcommand{\scaleMitMstLeveneP}{0.0005\xspace}                          % p-value (Levene) for a difference in spread of the errors between MIT and MST
\newcommand{\scaleMitMstAbsP}{0.020\xspace}                              % p-value (Mann-Whitney) for a difference in absolute error between MIT and MST
\newcommand{\scaleMitMstLeveneLightP}{0.007\xspace}                      % Levene p-value for MIT vs MST spread, mixes without the heaviest class only
\newcommand{\scaleModelEffectP}{0.0007\xspace}                           % ANOVA p-value of the model effect on the departure from additivity (controlling test and heavy class)
\newcommand{\scaleModelEffectF}{7.8\xspace}                              % ANOVA F statistic of the model effect on the departure from additivity
\newcommand{\scaleKindEffectP}{0.783\xspace}                             % ANOVA p-value of the test (MIT/MST) effect on the departure from additivity
\newcommand{\scaleModelEffectLightP}{0.544\xspace}                       % ANOVA p-value of the model effect, mixes without the heaviest class only
\newcommand{\scaleFamilyKwP}{0.849\xspace}                               % Kruskal-Wallis p-value for differences between mix-design families (equal shares, by cost, 4-5 classes)
\newcommand{\scaleHeavyDeepSeekMIT}{\ensuremath{-}12\,\%\xspace}         % mean departure from the harmonic law of DeepSeek-R1-Distill-Qwen-7B mix designs containing the heaviest class (600-1000 in, 1000-1500 out), MIT (2-minute instantaneous test)
\newcommand{\scaleHeavyDeepSeekMST}{\ensuremath{-}8\,\%\xspace}          % mean departure from the harmonic law of DeepSeek-R1-Distill-Qwen-7B mix designs containing the heaviest class (600-1000 in, 1000-1500 out), MST (30-minute sustainable test)
\newcommand{\scaleHeavyGemmaMIT}{+16\,\%\xspace}                         % mean departure from the harmonic law of Gemma-7B mix designs containing the heaviest class (600-1000 in, 1000-1500 out), MIT (2-minute instantaneous test)
\newcommand{\scaleHeavyGemmaMST}{+64\,\%\xspace}                         % mean departure from the harmonic law of Gemma-7B mix designs containing the heaviest class (600-1000 in, 1000-1500 out), MST (30-minute sustainable test)
\newcommand{\scaleHeavyLlamaMIT}{\ensuremath{-}9\,\%\xspace}             % mean departure from the harmonic law of Llama-3.1-8B-Instruct mix designs containing the heaviest class (600-1000 in, 1000-1500 out), MIT (2-minute instantaneous test)
\newcommand{\scaleHeavyLlamaMST}{\ensuremath{-}24\,\%\xspace}            % mean departure from the harmonic law of Llama-3.1-8B-Instruct mix designs containing the heaviest class (600-1000 in, 1000-1500 out), MST (30-minute sustainable test)
\newcommand{\scaleBurstMedian}{3.50\xspace}                              % median burst factor B = MIT/MST over all request classes and models (frequentist)
\newcommand{\scaleBurstQOne}{3.41\xspace}                                % first quartile of the burst factor over all classes
\newcommand{\scaleBurstQThree}{3.94\xspace}                              % third quartile of the burst factor over all classes
\newcommand{\scaleBurstOnBand}{77\,\%\xspace}                            % share of request classes whose B is within +-15 % of their model's median
\newcommand{\scaleBurstNCells}{60\xspace}                                % number of request classes (cells) for which B was measured (3 models x 20 classes)
\newcommand{\scaleBurstDeepSeek}{3.57\xspace}                            % median burst factor B = MIT/MST of DeepSeek-R1-Distill-Qwen-7B over its 20 classes (Bayesian when available)
\newcommand{\scaleBurstMaxDeepSeek}{5.6\xspace}                          % largest burst factor among DeepSeek-R1-Distill-Qwen-7B's classes
\newcommand{\scaleBurstOnBandDeepSeek}{70\,\%\xspace}                    % share of DeepSeek-R1-Distill-Qwen-7B's classes with B within +-15 % of its median
\newcommand{\scaleBurstDeepSeekLo}{3.49\xspace}                          % lower bound of the 95 % credible interval of DeepSeek-R1-Distill-Qwen-7B's median burst factor
\newcommand{\scaleBurstDeepSeekHi}{3.64\xspace}                          % upper bound of the 95 % credible interval of DeepSeek-R1-Distill-Qwen-7B's median burst factor
\newcommand{\scaleBurstGemma}{3.56\xspace}                               % median burst factor B = MIT/MST of Gemma-7B over its 20 classes (Bayesian when available)
\newcommand{\scaleBurstMaxGemma}{11.6\xspace}                            % largest burst factor among Gemma-7B's classes
\newcommand{\scaleBurstOnBandGemma}{75\,\%\xspace}                       % share of Gemma-7B's classes with B within +-15 % of its median
\newcommand{\scaleBurstGemmaLo}{3.49\xspace}                             % lower bound of the 95 % credible interval of Gemma-7B's median burst factor
\newcommand{\scaleBurstGemmaHi}{3.65\xspace}                             % upper bound of the 95 % credible interval of Gemma-7B's median burst factor
\newcommand{\scaleBurstLlama}{3.55\xspace}                               % median burst factor B = MIT/MST of Llama-3.1-8B-Instruct over its 20 classes (Bayesian when available)
\newcommand{\scaleBurstMaxLlama}{5.4\xspace}                             % largest burst factor among Llama-3.1-8B-Instruct's classes
\newcommand{\scaleBurstOnBandLlama}{85\,\%\xspace}                       % share of Llama-3.1-8B-Instruct's classes with B within +-15 % of its median
\newcommand{\scaleBurstLlamaLo}{3.48\xspace}                             % lower bound of the 95 % credible interval of Llama-3.1-8B-Instruct's median burst factor
\newcommand{\scaleBurstLlamaHi}{3.62\xspace}                             % upper bound of the 95 % credible interval of Llama-3.1-8B-Instruct's median burst factor
\newcommand{\scaleBurstModelP}{0.742\xspace}                             % ANOVA p-value of the model effect on log B (controlling input and output range)
\newcommand{\scaleBurstOutP}{0.139\xspace}                               % ANOVA p-value of the output-length effect on log B
\newcommand{\scaleBurstInP}{0.672\xspace}                                % ANOVA p-value of the input-length effect on log B
\newcommand{\scaleBurstOutRobustP}{\ensuremath{6\times 10^{-7}}\xspace}  % ANOVA p-value of the output-length effect on log B without Gemma's two anomalous cells
\newcommand{\scaleBurstModelRobustP}{0.045\xspace}                       % ANOVA p-value of the model effect on log B without Gemma's two anomalous cells
\newcommand{\scaleBurstShort}{4.15\xspace}                               % geometric-mean burst factor of classes with short answers (1-100 output tokens)
\newcommand{\scaleBurstRest}{3.68\xspace}                                % geometric-mean burst factor of the other classes
\newcommand{\scaleBurstShortRatio}{+13\,\%\xspace}                       % how much larger B is for short answers than for the rest
\newcommand{\scaleBurstShortP}{0.0004\xspace}                            % Mann-Whitney p-value for short answers vs the rest
\newcommand{\scaleBurstPaperExample}{1.4\xspace}                         % burst factor implied by the paper's introductory example (1,400 req/min vs 30,000 req/30 min)
\newcommand{\scaleBurstMixN}{12\xspace}                                  % number of mix designs measured with the same composition in both MIT and MST
\newcommand{\scaleBurstMixErr}{+19\,\%\xspace}                           % mean excess of the measured burst factor of mixes over the one predicted from their classes
\newcommand{\scaleBurstMixWithin}{50\,\%\xspace}                         % share of those mixes whose measured B is within +-15 % of the prediction
\newcommand{\scaleBurstEquivMin}{100\,\%\xspace}                         % smallest posterior probability, over model pairs, that two models' burst factors are equivalent within +-10 %
\newcommand{\scaleBurstPairMaxDiff}{0.5\,\%\xspace}                      % largest posterior-median difference between the burst factors of two models
\newcommand{\scaleBurstPairCi}{3.5\,\%\xspace}                           % largest distance from 1 of the 95 % credible bounds of the ratio between two models' burst factors
\newcommand{\scaleTbWindow}{2\xspace}                                    % MIT window W in minutes used to characterise bursts
\newcommand{\scaleTbDepth}{5.0\xspace}                                   % token-bucket depth b/s in minutes of refill, b = (B-1) W, with the typical B
\newcommand{\scaleTbDepthSafe}{4.6\xspace}                               % token-bucket depth in minutes of refill with the conservative B_safe
\newcommand{\scaleTbDepthSmall}{6.1\xspace}                              % token-bucket depth in minutes of refill for plans of requests up to 300 tokens
\newcommand{\scaleTbBSafe}{3.29\xspace}                                  % conservative burst factor B_safe: 10th percentile over classes (covers 90 % of them)
\newcommand{\scaleTbBSmall}{4.05\xspace}                                 % median burst factor of classes up to 300 input and output tokens (the paper's plans)
\newcommand{\scaleTbRho}{1.2\xspace}                                     % oversubscription ratio rho used in the paper's example
\newcommand{\scaleTbPhiMax}{77\,\%\xspace}                               % largest share of the allocation that may burst at once at rho = TbRho: phi <= (B/rho - 1)/(B - 1)
\newcommand{\scaleTbRhoMaxHalf}{1.56\xspace}                             % largest burst-safe oversubscription ratio if half of the allocation bursts at once: rho <= B/(1 + phi(B-1))
\newcommand{\scaleTbRhoMaxQuarter}{2.15\xspace}                          % largest burst-safe oversubscription ratio if a quarter of the allocation bursts at once
\newcommand{\scaleTbPaperRatio}{21\xspace}                               % burst-to-sustained rate ratio implied by Table 3 of the paper (Basic plan)
\newcommand{\scaleTbSkMaxDeepSeek}{2{,}307\xspace}                       % SK_max of DeepSeek-R1-Distill-Qwen-7B: largest MST in its capacity matrix, req/min (1 K per request of that class)
\newcommand{\scaleTbThrottleDeepSeek}{8{,}226\xspace}                    % throttling threshold Theta = B * SK_max of DeepSeek-R1-Distill-Qwen-7B, sustainable units per minute averaged over W
\newcommand{\scaleTbSkMaxGemma}{2{,}047\xspace}                          % SK_max of Gemma-7B: largest MST in its capacity matrix, req/min (1 K per request of that class)
\newcommand{\scaleTbThrottleGemma}{7{,}297\xspace}                       % throttling threshold Theta = B * SK_max of Gemma-7B, sustainable units per minute averaged over W
\newcommand{\scaleTbSkMaxLlama}{2{,}176\xspace}                          % SK_max of Llama-3.1-8B-Instruct: largest MST in its capacity matrix, req/min (1 K per request of that class)
\newcommand{\scaleTbThrottleLlama}{7{,}717\xspace}                       % throttling threshold Theta = B * SK_max of Llama-3.1-8B-Instruct, sustainable units per minute averaged over W
\newcommand{\scaleTbBasicBurstTwo}{14\xspace}                            % burst allowance of the paper's Basic plan (50 req/30 min) derived from B, in requests per 2 minutes
\newcommand{\scaleTbBasicDepth}{10\xspace}                               % token-bucket depth of the Basic plan derived from B, in requests
\newcommand{\scaleTbProBurstTwo}{47\xspace}                              % burst allowance of the paper's Pro plan (175 req/30 min) derived from B, in requests per 2 minutes
\newcommand{\scaleTbProDepth}{36\xspace}                                 % token-bucket depth of the Pro plan derived from B, in requests
 in the preamble; every macro is prefixed with \scale.
% Each line ends with a comment saying what the value means, measures or represents
% (descriptions live in analysis/macro_docs.py). Percentages are relative differences
% unless stated; 'design level' means one value per mix design, repeats averaged.
\providecommand{\xspace}{}
\newcommand{\scaleNEvaluations}{7{,}944\xspace}                          % number of MoST pass/fail SLA evaluations analysed (all runs, after removing exact copies)
\newcommand{\scaleNMixRuns}{301\xspace}                                  % number of mix runs (MoST searches on request mixes)
\newcommand{\scaleNDesigns}{114\xspace}                                  % number of distinct mix designs (model x test x composition); unit of the frequentist tests
\newcommand{\scaleNRepeatRuns}{270\xspace}                               % number of mix runs that belong to designs measured more than once
\newcommand{\scaleNRepeatDesigns}{83\xspace}                             % number of mix designs measured more than once (stability check)
\newcommand{\scaleNDuplicates}{114\xspace}                               % number of byte-identical copies of a run found in the archive and counted once
\newcommand{\scaleNPureRuns}{360\xspace}                                 % number of pure-class runs (one request class per MoST search)
\newcommand{\scaleQHatFreq}{\ensuremath{-}0.91\xspace}                   % power-mean exponent q of the combination law, profile-likelihood estimate (q=-1 harmonic, 0 geometric, +1 arithmetic)
\newcommand{\scaleQFreqLo}{\ensuremath{-}0.98\xspace}                    % lower bound of the 95 % profile-likelihood confidence interval of q
\newcommand{\scaleQFreqHi}{\ensuremath{-}0.84\xspace}                    % upper bound of the 95 % profile-likelihood confidence interval of q
\newcommand{\scaleHarmonicWinShare}{97\,\%\xspace}                       % share of mix runs for which the harmonic prediction is closer to the measurement than the arithmetic one
\newcommand{\scaleCalibSlope}{0.96\xspace}                               % slope of the regression log(measured mix capacity) ~ log(harmonic prediction); 1 = correct scaling
\newcommand{\scaleCalibSlopeLo}{0.90\xspace}                             % lower bound of the 95 % confidence interval of the calibration slope (cluster-robust by model)
\newcommand{\scaleCalibSlopeHi}{1.02\xspace}                             % upper bound of the 95 % confidence interval of the calibration slope
\newcommand{\scaleCalibRsq}{0.97\xspace}                                 % R^2 of the calibration regression on the log scale
\newcommand{\scaleBiasFreq}{\ensuremath{-}3.0\,\%\xspace}                % mean departure of measured mix capacity from the harmonic prediction, frequentist, designs as units
\newcommand{\scaleBiasFreqLo}{\ensuremath{-}6.9\,\%\xspace}              % lower bound of the 95 % confidence interval of BiasFreq
\newcommand{\scaleBiasFreqHi}{+1.0\,\%\xspace}                           % upper bound of the 95 % confidence interval of BiasFreq
\newcommand{\scaleBiasTostP}{0.001\xspace}                               % p-value of the TOST equivalence test of the bias against +-10 % (small = practically zero)
\newcommand{\scaleArithMeanErr}{\ensuremath{-}53\,\%\xspace}             % mean error of the arithmetic-mean law (measured vs predicted mix capacity, design level); negative = overestimates
\newcommand{\scaleGeomMeanErr}{\ensuremath{-}39\,\%\xspace}              % mean error of the geometric-mean law (design level); negative = overestimates
\newcommand{\scaleBoundsIn}{294\xspace}                                  % number of mix runs whose capacity lies between their slowest and fastest class
\newcommand{\scaleBoundsN}{301\xspace}                                   % number of mix runs checked for boundedness
\newcommand{\scaleBoundsShare}{97.7\,\%\xspace}                          % share of mix runs within the capacity range of their classes
\newcommand{\scaleAboveArith}{0\xspace}                                  % number of mix runs whose capacity exceeds the arithmetic mean of their classes
\newcommand{\scaleQHatBayes}{\ensuremath{-}0.96\xspace}                  % power-mean exponent q, Bayesian posterior median (free-q model)
\newcommand{\scaleQBayesLo}{\ensuremath{-}1.02\xspace}                   % lower bound of the 95 % credible interval of q
\newcommand{\scaleQBayesHi}{\ensuremath{-}0.88\xspace}                   % upper bound of the 95 % credible interval of q
\newcommand{\scaleBiasBayes}{\ensuremath{-}5.5\,\%\xspace}               % pooled bias delta_0 of mixes w.r.t. the harmonic law, Bayesian posterior median
\newcommand{\scaleBiasBayesLo}{\ensuremath{-}9.9\,\%\xspace}             % lower bound of the 90 % credible interval of delta_0
\newcommand{\scaleBiasBayesHi}{\ensuremath{-}1.0\,\%\xspace}             % upper bound of the 90 % credible interval of delta_0
\newcommand{\scaleBiasBayesPRope}{92\,\%\xspace}                         % posterior probability that |delta_0| < 10 % (region of practical equivalence)
\newcommand{\scalePredFree}{+1.7\xspace}                                 % out-of-sample log predictive density of the free-q law minus that of the harmonic law (higher = better)
\newcommand{\scalePredFreeSE}{3.3\xspace}                                % standard error of PredFree
\newcommand{\scalePredGeom}{\ensuremath{-}264.5\xspace}                  % out-of-sample log predictive density of the geometric law minus that of the harmonic law
\newcommand{\scalePredGeomSE}{24.8\xspace}                               % standard error of PredGeom
\newcommand{\scalePredArith}{\ensuremath{-}441.4\xspace}                 % out-of-sample log predictive density of the arithmetic law minus that of the harmonic law
\newcommand{\scalePredArithSE}{38.1\xspace}                              % standard error of PredArith
\newcommand{\scaleSdMixMIT}{6\,\%\xspace}                                % posterior SD of the design effect (departure of a mix design from the harmonic law), MIT (2-minute instantaneous test), as %
\newcommand{\scaleSdRunMIT}{0.3\,\%\xspace}                              % posterior SD between repeated runs of the same mix design, MIT (2-minute instantaneous test), as %
\newcommand{\scaleSdRepMIT}{0.3\,\%\xspace}                              % posterior SD between repeated runs of the same pure class, MIT (2-minute instantaneous test), as %
\newcommand{\scaleBandMIT}{\ensuremath{\pm}15\,\%\xspace}                % half-width of the 90 % planning band around the harmonic prediction, MIT (2-minute instantaneous test)
\newcommand{\scaleSdMixMST}{21\,\%\xspace}                               % posterior SD of the design effect (departure of a mix design from the harmonic law), MST (30-minute sustainable test), as %
\newcommand{\scaleSdRunMST}{1.8\,\%\xspace}                              % posterior SD between repeated runs of the same mix design, MST (30-minute sustainable test), as %
\newcommand{\scaleSdRepMST}{1.7\,\%\xspace}                              % posterior SD between repeated runs of the same pure class, MST (30-minute sustainable test), as %
\newcommand{\scaleBandMST}{\ensuremath{\pm}39\,\%\xspace}                % half-width of the 90 % planning band around the harmonic prediction, MST (30-minute sustainable test)
\newcommand{\scaleBiasDeepSeekMIT}{\ensuremath{-}8.8\,\%\xspace}         % Bayesian bias delta of DeepSeek-R1-Distill-Qwen-7B mixes w.r.t. the harmonic law, MIT (2-minute instantaneous test)
\newcommand{\scaleBiasDeepSeekMST}{\ensuremath{-}4.1\,\%\xspace}         % Bayesian bias delta of DeepSeek-R1-Distill-Qwen-7B mixes w.r.t. the harmonic law, MST (30-minute sustainable test)
\newcommand{\scaleBiasGemmaMIT}{\ensuremath{-}0.3\,\%\xspace}            % Bayesian bias delta of Gemma-7B mixes w.r.t. the harmonic law, MIT (2-minute instantaneous test)
\newcommand{\scaleBiasGemmaMST}{\ensuremath{-}4.6\,\%\xspace}            % Bayesian bias delta of Gemma-7B mixes w.r.t. the harmonic law, MST (30-minute sustainable test)
\newcommand{\scaleBiasLlamaMIT}{\ensuremath{-}7.0\,\%\xspace}            % Bayesian bias delta of Llama-3.1-8B-Instruct mixes w.r.t. the harmonic law, MIT (2-minute instantaneous test)
\newcommand{\scaleBiasLlamaMST}{\ensuremath{-}7.8\,\%\xspace}            % Bayesian bias delta of Llama-3.1-8B-Instruct mixes w.r.t. the harmonic law, MST (30-minute sustainable test)
\newcommand{\scaleBayesDraws}{600\xspace}                                % NUTS draws per chain (4 chains, same number of tuning steps)
\newcommand{\scaleMaxRhat}{1.02\xspace}                                  % largest split-R-hat over the global parameters of the reported Bayesian models (<1.05 = converged)
\newcommand{\scaleMaeMIT}{10\,\%\xspace}                                 % mean absolute error of the harmonic prediction per mix design, MIT (2-minute instantaneous test)
\newcommand{\scaleBiasDesignMIT}{\ensuremath{-}2.9\,\%\xspace}           % mean departure from the harmonic law per mix design, MIT (2-minute instantaneous test)
\newcommand{\scaleWithinTenMIT}{61\,\%\xspace}                           % share of mix designs predicted within +-10 %, MIT (2-minute instantaneous test)
\newcommand{\scaleWithinTenMITLight}{75\,\%\xspace}                      % share of mix designs without the heaviest class predicted within +-10 %, MIT (2-minute instantaneous test)
\newcommand{\scaleWithinTenMITHeavy}{42\,\%\xspace}                      % share of mix designs with the heaviest class predicted within +-10 %, MIT (2-minute instantaneous test)
\newcommand{\scaleMaeMST}{23\,\%\xspace}                                 % mean absolute error of the harmonic prediction per mix design, MST (30-minute sustainable test)
\newcommand{\scaleBiasDesignMST}{\ensuremath{-}3.1\,\%\xspace}           % mean departure from the harmonic law per mix design, MST (30-minute sustainable test)
\newcommand{\scaleWithinTenMST}{48\,\%\xspace}                           % share of mix designs predicted within +-10 %, MST (30-minute sustainable test)
\newcommand{\scaleWithinTenMSTLight}{58\,\%\xspace}                      % share of mix designs without the heaviest class predicted within +-10 %, MST (30-minute sustainable test)
\newcommand{\scaleWithinTenMSTHeavy}{25\,\%\xspace}                      % share of mix designs with the heaviest class predicted within +-10 %, MST (30-minute sustainable test)
\newcommand{\scaleMitMstMeanP}{0.507\xspace}                             % p-value (Mann-Whitney) for a difference in mean bias between MIT and MST designs
\newcommand{\scaleMitMstLeveneP}{0.0005\xspace}                          % p-value (Levene) for a difference in spread of the errors between MIT and MST
\newcommand{\scaleMitMstAbsP}{0.020\xspace}                              % p-value (Mann-Whitney) for a difference in absolute error between MIT and MST
\newcommand{\scaleMitMstLeveneLightP}{0.007\xspace}                      % Levene p-value for MIT vs MST spread, mixes without the heaviest class only
\newcommand{\scaleModelEffectP}{0.0007\xspace}                           % ANOVA p-value of the model effect on the departure from additivity (controlling test and heavy class)
\newcommand{\scaleModelEffectF}{7.8\xspace}                              % ANOVA F statistic of the model effect on the departure from additivity
\newcommand{\scaleKindEffectP}{0.783\xspace}                             % ANOVA p-value of the test (MIT/MST) effect on the departure from additivity
\newcommand{\scaleModelEffectLightP}{0.544\xspace}                       % ANOVA p-value of the model effect, mixes without the heaviest class only
\newcommand{\scaleFamilyKwP}{0.849\xspace}                               % Kruskal-Wallis p-value for differences between mix-design families (equal shares, by cost, 4-5 classes)
\newcommand{\scaleHeavyDeepSeekMIT}{\ensuremath{-}12\,\%\xspace}         % mean departure from the harmonic law of DeepSeek-R1-Distill-Qwen-7B mix designs containing the heaviest class (600-1000 in, 1000-1500 out), MIT (2-minute instantaneous test)
\newcommand{\scaleHeavyDeepSeekMST}{\ensuremath{-}8\,\%\xspace}          % mean departure from the harmonic law of DeepSeek-R1-Distill-Qwen-7B mix designs containing the heaviest class (600-1000 in, 1000-1500 out), MST (30-minute sustainable test)
\newcommand{\scaleHeavyGemmaMIT}{+16\,\%\xspace}                         % mean departure from the harmonic law of Gemma-7B mix designs containing the heaviest class (600-1000 in, 1000-1500 out), MIT (2-minute instantaneous test)
\newcommand{\scaleHeavyGemmaMST}{+64\,\%\xspace}                         % mean departure from the harmonic law of Gemma-7B mix designs containing the heaviest class (600-1000 in, 1000-1500 out), MST (30-minute sustainable test)
\newcommand{\scaleHeavyLlamaMIT}{\ensuremath{-}9\,\%\xspace}             % mean departure from the harmonic law of Llama-3.1-8B-Instruct mix designs containing the heaviest class (600-1000 in, 1000-1500 out), MIT (2-minute instantaneous test)
\newcommand{\scaleHeavyLlamaMST}{\ensuremath{-}24\,\%\xspace}            % mean departure from the harmonic law of Llama-3.1-8B-Instruct mix designs containing the heaviest class (600-1000 in, 1000-1500 out), MST (30-minute sustainable test)
\newcommand{\scaleBurstMedian}{3.50\xspace}                              % median burst factor B = MIT/MST over all request classes and models (frequentist)
\newcommand{\scaleBurstQOne}{3.41\xspace}                                % first quartile of the burst factor over all classes
\newcommand{\scaleBurstQThree}{3.94\xspace}                              % third quartile of the burst factor over all classes
\newcommand{\scaleBurstOnBand}{77\,\%\xspace}                            % share of request classes whose B is within +-15 % of their model's median
\newcommand{\scaleBurstNCells}{60\xspace}                                % number of request classes (cells) for which B was measured (3 models x 20 classes)
\newcommand{\scaleBurstDeepSeek}{3.57\xspace}                            % median burst factor B = MIT/MST of DeepSeek-R1-Distill-Qwen-7B over its 20 classes (Bayesian when available)
\newcommand{\scaleBurstMaxDeepSeek}{5.6\xspace}                          % largest burst factor among DeepSeek-R1-Distill-Qwen-7B's classes
\newcommand{\scaleBurstOnBandDeepSeek}{70\,\%\xspace}                    % share of DeepSeek-R1-Distill-Qwen-7B's classes with B within +-15 % of its median
\newcommand{\scaleBurstDeepSeekLo}{3.49\xspace}                          % lower bound of the 95 % credible interval of DeepSeek-R1-Distill-Qwen-7B's median burst factor
\newcommand{\scaleBurstDeepSeekHi}{3.64\xspace}                          % upper bound of the 95 % credible interval of DeepSeek-R1-Distill-Qwen-7B's median burst factor
\newcommand{\scaleBurstGemma}{3.56\xspace}                               % median burst factor B = MIT/MST of Gemma-7B over its 20 classes (Bayesian when available)
\newcommand{\scaleBurstMaxGemma}{11.6\xspace}                            % largest burst factor among Gemma-7B's classes
\newcommand{\scaleBurstOnBandGemma}{75\,\%\xspace}                       % share of Gemma-7B's classes with B within +-15 % of its median
\newcommand{\scaleBurstGemmaLo}{3.49\xspace}                             % lower bound of the 95 % credible interval of Gemma-7B's median burst factor
\newcommand{\scaleBurstGemmaHi}{3.65\xspace}                             % upper bound of the 95 % credible interval of Gemma-7B's median burst factor
\newcommand{\scaleBurstLlama}{3.55\xspace}                               % median burst factor B = MIT/MST of Llama-3.1-8B-Instruct over its 20 classes (Bayesian when available)
\newcommand{\scaleBurstMaxLlama}{5.4\xspace}                             % largest burst factor among Llama-3.1-8B-Instruct's classes
\newcommand{\scaleBurstOnBandLlama}{85\,\%\xspace}                       % share of Llama-3.1-8B-Instruct's classes with B within +-15 % of its median
\newcommand{\scaleBurstLlamaLo}{3.48\xspace}                             % lower bound of the 95 % credible interval of Llama-3.1-8B-Instruct's median burst factor
\newcommand{\scaleBurstLlamaHi}{3.62\xspace}                             % upper bound of the 95 % credible interval of Llama-3.1-8B-Instruct's median burst factor
\newcommand{\scaleBurstModelP}{0.742\xspace}                             % ANOVA p-value of the model effect on log B (controlling input and output range)
\newcommand{\scaleBurstOutP}{0.139\xspace}                               % ANOVA p-value of the output-length effect on log B
\newcommand{\scaleBurstInP}{0.672\xspace}                                % ANOVA p-value of the input-length effect on log B
\newcommand{\scaleBurstOutRobustP}{\ensuremath{6\times 10^{-7}}\xspace}  % ANOVA p-value of the output-length effect on log B without Gemma's two anomalous cells
\newcommand{\scaleBurstModelRobustP}{0.045\xspace}                       % ANOVA p-value of the model effect on log B without Gemma's two anomalous cells
\newcommand{\scaleBurstShort}{4.15\xspace}                               % geometric-mean burst factor of classes with short answers (1-100 output tokens)
\newcommand{\scaleBurstRest}{3.68\xspace}                                % geometric-mean burst factor of the other classes
\newcommand{\scaleBurstShortRatio}{+13\,\%\xspace}                       % how much larger B is for short answers than for the rest
\newcommand{\scaleBurstShortP}{0.0004\xspace}                            % Mann-Whitney p-value for short answers vs the rest
\newcommand{\scaleBurstPaperExample}{1.4\xspace}                         % burst factor implied by the paper's introductory example (1,400 req/min vs 30,000 req/30 min)
\newcommand{\scaleBurstMixN}{12\xspace}                                  % number of mix designs measured with the same composition in both MIT and MST
\newcommand{\scaleBurstMixErr}{+19\,\%\xspace}                           % mean excess of the measured burst factor of mixes over the one predicted from their classes
\newcommand{\scaleBurstMixWithin}{50\,\%\xspace}                         % share of those mixes whose measured B is within +-15 % of the prediction
\newcommand{\scaleBurstEquivMin}{100\,\%\xspace}                         % smallest posterior probability, over model pairs, that two models' burst factors are equivalent within +-10 %
\newcommand{\scaleBurstPairMaxDiff}{0.5\,\%\xspace}                      % largest posterior-median difference between the burst factors of two models
\newcommand{\scaleBurstPairCi}{3.5\,\%\xspace}                           % largest distance from 1 of the 95 % credible bounds of the ratio between two models' burst factors
\newcommand{\scaleTbWindow}{2\xspace}                                    % MIT window W in minutes used to characterise bursts
\newcommand{\scaleTbDepth}{5.0\xspace}                                   % token-bucket depth b/s in minutes of refill, b = (B-1) W, with the typical B
\newcommand{\scaleTbDepthSafe}{4.6\xspace}                               % token-bucket depth in minutes of refill with the conservative B_safe
\newcommand{\scaleTbDepthSmall}{6.1\xspace}                              % token-bucket depth in minutes of refill for plans of requests up to 300 tokens
\newcommand{\scaleTbBSafe}{3.29\xspace}                                  % conservative burst factor B_safe: 10th percentile over classes (covers 90 % of them)
\newcommand{\scaleTbBSmall}{4.05\xspace}                                 % median burst factor of classes up to 300 input and output tokens (the paper's plans)
\newcommand{\scaleTbRho}{1.2\xspace}                                     % oversubscription ratio rho used in the paper's example
\newcommand{\scaleTbPhiMax}{77\,\%\xspace}                               % largest share of the allocation that may burst at once at rho = TbRho: phi <= (B/rho - 1)/(B - 1)
\newcommand{\scaleTbRhoMaxHalf}{1.56\xspace}                             % largest burst-safe oversubscription ratio if half of the allocation bursts at once: rho <= B/(1 + phi(B-1))
\newcommand{\scaleTbRhoMaxQuarter}{2.15\xspace}                          % largest burst-safe oversubscription ratio if a quarter of the allocation bursts at once
\newcommand{\scaleTbPaperRatio}{21\xspace}                               % burst-to-sustained rate ratio implied by Table 3 of the paper (Basic plan)
\newcommand{\scaleTbSkMaxDeepSeek}{2{,}307\xspace}                       % SK_max of DeepSeek-R1-Distill-Qwen-7B: largest MST in its capacity matrix, req/min (1 K per request of that class)
\newcommand{\scaleTbThrottleDeepSeek}{8{,}226\xspace}                    % throttling threshold Theta = B * SK_max of DeepSeek-R1-Distill-Qwen-7B, sustainable units per minute averaged over W
\newcommand{\scaleTbSkMaxGemma}{2{,}047\xspace}                          % SK_max of Gemma-7B: largest MST in its capacity matrix, req/min (1 K per request of that class)
\newcommand{\scaleTbThrottleGemma}{7{,}297\xspace}                       % throttling threshold Theta = B * SK_max of Gemma-7B, sustainable units per minute averaged over W
\newcommand{\scaleTbSkMaxLlama}{2{,}176\xspace}                          % SK_max of Llama-3.1-8B-Instruct: largest MST in its capacity matrix, req/min (1 K per request of that class)
\newcommand{\scaleTbThrottleLlama}{7{,}717\xspace}                       % throttling threshold Theta = B * SK_max of Llama-3.1-8B-Instruct, sustainable units per minute averaged over W
\newcommand{\scaleTbBasicBurstTwo}{14\xspace}                            % burst allowance of the paper's Basic plan (50 req/30 min) derived from B, in requests per 2 minutes
\newcommand{\scaleTbBasicDepth}{10\xspace}                               % token-bucket depth of the Basic plan derived from B, in requests
\newcommand{\scaleTbProBurstTwo}{47\xspace}                              % burst allowance of the paper's Pro plan (175 req/30 min) derived from B, in requests per 2 minutes
\newcommand{\scaleTbProDepth}{36\xspace}                                 % token-bucket depth of the Pro plan derived from B, in requests
 in the preamble; every macro is prefixed with \scale.
% Each line ends with a comment saying what the value means, measures or represents
% (descriptions live in analysis/macro_docs.py). Percentages are relative differences
% unless stated; 'design level' means one value per mix design, repeats averaged.
\providecommand{\xspace}{}
\newcommand{\scaleNEvaluations}{7{,}944\xspace}                          % number of MoST pass/fail SLA evaluations analysed (all runs, after removing exact copies)
\newcommand{\scaleNMixRuns}{301\xspace}                                  % number of mix runs (MoST searches on request mixes)
\newcommand{\scaleNDesigns}{114\xspace}                                  % number of distinct mix designs (model x test x composition); unit of the frequentist tests
\newcommand{\scaleNRepeatRuns}{270\xspace}                               % number of mix runs that belong to designs measured more than once
\newcommand{\scaleNRepeatDesigns}{83\xspace}                             % number of mix designs measured more than once (stability check)
\newcommand{\scaleNDuplicates}{114\xspace}                               % number of byte-identical copies of a run found in the archive and counted once
\newcommand{\scaleNPureRuns}{360\xspace}                                 % number of pure-class runs (one request class per MoST search)
\newcommand{\scaleQHatFreq}{\ensuremath{-}0.91\xspace}                   % power-mean exponent q of the combination law, profile-likelihood estimate (q=-1 harmonic, 0 geometric, +1 arithmetic)
\newcommand{\scaleQFreqLo}{\ensuremath{-}0.98\xspace}                    % lower bound of the 95 % profile-likelihood confidence interval of q
\newcommand{\scaleQFreqHi}{\ensuremath{-}0.84\xspace}                    % upper bound of the 95 % profile-likelihood confidence interval of q
\newcommand{\scaleHarmonicWinShare}{97\,\%\xspace}                       % share of mix runs for which the harmonic prediction is closer to the measurement than the arithmetic one
\newcommand{\scaleCalibSlope}{0.96\xspace}                               % slope of the regression log(measured mix capacity) ~ log(harmonic prediction); 1 = correct scaling
\newcommand{\scaleCalibSlopeLo}{0.90\xspace}                             % lower bound of the 95 % confidence interval of the calibration slope (cluster-robust by model)
\newcommand{\scaleCalibSlopeHi}{1.02\xspace}                             % upper bound of the 95 % confidence interval of the calibration slope
\newcommand{\scaleCalibRsq}{0.97\xspace}                                 % R^2 of the calibration regression on the log scale
\newcommand{\scaleBiasFreq}{\ensuremath{-}3.0\,\%\xspace}                % mean departure of measured mix capacity from the harmonic prediction, frequentist, designs as units
\newcommand{\scaleBiasFreqLo}{\ensuremath{-}6.9\,\%\xspace}              % lower bound of the 95 % confidence interval of BiasFreq
\newcommand{\scaleBiasFreqHi}{+1.0\,\%\xspace}                           % upper bound of the 95 % confidence interval of BiasFreq
\newcommand{\scaleBiasTostP}{0.001\xspace}                               % p-value of the TOST equivalence test of the bias against +-10 % (small = practically zero)
\newcommand{\scaleArithMeanErr}{\ensuremath{-}53\,\%\xspace}             % mean error of the arithmetic-mean law (measured vs predicted mix capacity, design level); negative = overestimates
\newcommand{\scaleGeomMeanErr}{\ensuremath{-}39\,\%\xspace}              % mean error of the geometric-mean law (design level); negative = overestimates
\newcommand{\scaleBoundsIn}{294\xspace}                                  % number of mix runs whose capacity lies between their slowest and fastest class
\newcommand{\scaleBoundsN}{301\xspace}                                   % number of mix runs checked for boundedness
\newcommand{\scaleBoundsShare}{97.7\,\%\xspace}                          % share of mix runs within the capacity range of their classes
\newcommand{\scaleAboveArith}{0\xspace}                                  % number of mix runs whose capacity exceeds the arithmetic mean of their classes
\newcommand{\scaleQHatBayes}{\ensuremath{-}0.96\xspace}                  % power-mean exponent q, Bayesian posterior median (free-q model)
\newcommand{\scaleQBayesLo}{\ensuremath{-}1.02\xspace}                   % lower bound of the 95 % credible interval of q
\newcommand{\scaleQBayesHi}{\ensuremath{-}0.88\xspace}                   % upper bound of the 95 % credible interval of q
\newcommand{\scaleBiasBayes}{\ensuremath{-}5.5\,\%\xspace}               % pooled bias delta_0 of mixes w.r.t. the harmonic law, Bayesian posterior median
\newcommand{\scaleBiasBayesLo}{\ensuremath{-}9.9\,\%\xspace}             % lower bound of the 90 % credible interval of delta_0
\newcommand{\scaleBiasBayesHi}{\ensuremath{-}1.0\,\%\xspace}             % upper bound of the 90 % credible interval of delta_0
\newcommand{\scaleBiasBayesPRope}{92\,\%\xspace}                         % posterior probability that |delta_0| < 10 % (region of practical equivalence)
\newcommand{\scalePredFree}{+1.7\xspace}                                 % out-of-sample log predictive density of the free-q law minus that of the harmonic law (higher = better)
\newcommand{\scalePredFreeSE}{3.3\xspace}                                % standard error of PredFree
\newcommand{\scalePredGeom}{\ensuremath{-}264.5\xspace}                  % out-of-sample log predictive density of the geometric law minus that of the harmonic law
\newcommand{\scalePredGeomSE}{24.8\xspace}                               % standard error of PredGeom
\newcommand{\scalePredArith}{\ensuremath{-}441.4\xspace}                 % out-of-sample log predictive density of the arithmetic law minus that of the harmonic law
\newcommand{\scalePredArithSE}{38.1\xspace}                              % standard error of PredArith
\newcommand{\scaleSdMixMIT}{6\,\%\xspace}                                % posterior SD of the design effect (departure of a mix design from the harmonic law), MIT (2-minute instantaneous test), as %
\newcommand{\scaleSdRunMIT}{0.3\,\%\xspace}                              % posterior SD between repeated runs of the same mix design, MIT (2-minute instantaneous test), as %
\newcommand{\scaleSdRepMIT}{0.3\,\%\xspace}                              % posterior SD between repeated runs of the same pure class, MIT (2-minute instantaneous test), as %
\newcommand{\scaleBandMIT}{\ensuremath{\pm}15\,\%\xspace}                % half-width of the 90 % planning band around the harmonic prediction, MIT (2-minute instantaneous test)
\newcommand{\scaleSdMixMST}{21\,\%\xspace}                               % posterior SD of the design effect (departure of a mix design from the harmonic law), MST (30-minute sustainable test), as %
\newcommand{\scaleSdRunMST}{1.8\,\%\xspace}                              % posterior SD between repeated runs of the same mix design, MST (30-minute sustainable test), as %
\newcommand{\scaleSdRepMST}{1.7\,\%\xspace}                              % posterior SD between repeated runs of the same pure class, MST (30-minute sustainable test), as %
\newcommand{\scaleBandMST}{\ensuremath{\pm}39\,\%\xspace}                % half-width of the 90 % planning band around the harmonic prediction, MST (30-minute sustainable test)
\newcommand{\scaleBiasDeepSeekMIT}{\ensuremath{-}8.8\,\%\xspace}         % Bayesian bias delta of DeepSeek-R1-Distill-Qwen-7B mixes w.r.t. the harmonic law, MIT (2-minute instantaneous test)
\newcommand{\scaleBiasDeepSeekMST}{\ensuremath{-}4.1\,\%\xspace}         % Bayesian bias delta of DeepSeek-R1-Distill-Qwen-7B mixes w.r.t. the harmonic law, MST (30-minute sustainable test)
\newcommand{\scaleBiasGemmaMIT}{\ensuremath{-}0.3\,\%\xspace}            % Bayesian bias delta of Gemma-7B mixes w.r.t. the harmonic law, MIT (2-minute instantaneous test)
\newcommand{\scaleBiasGemmaMST}{\ensuremath{-}4.6\,\%\xspace}            % Bayesian bias delta of Gemma-7B mixes w.r.t. the harmonic law, MST (30-minute sustainable test)
\newcommand{\scaleBiasLlamaMIT}{\ensuremath{-}7.0\,\%\xspace}            % Bayesian bias delta of Llama-3.1-8B-Instruct mixes w.r.t. the harmonic law, MIT (2-minute instantaneous test)
\newcommand{\scaleBiasLlamaMST}{\ensuremath{-}7.8\,\%\xspace}            % Bayesian bias delta of Llama-3.1-8B-Instruct mixes w.r.t. the harmonic law, MST (30-minute sustainable test)
\newcommand{\scaleBayesDraws}{600\xspace}                                % NUTS draws per chain (4 chains, same number of tuning steps)
\newcommand{\scaleMaxRhat}{1.02\xspace}                                  % largest split-R-hat over the global parameters of the reported Bayesian models (<1.05 = converged)
\newcommand{\scaleMaeMIT}{10\,\%\xspace}                                 % mean absolute error of the harmonic prediction per mix design, MIT (2-minute instantaneous test)
\newcommand{\scaleBiasDesignMIT}{\ensuremath{-}2.9\,\%\xspace}           % mean departure from the harmonic law per mix design, MIT (2-minute instantaneous test)
\newcommand{\scaleWithinTenMIT}{61\,\%\xspace}                           % share of mix designs predicted within +-10 %, MIT (2-minute instantaneous test)
\newcommand{\scaleWithinTenMITLight}{75\,\%\xspace}                      % share of mix designs without the heaviest class predicted within +-10 %, MIT (2-minute instantaneous test)
\newcommand{\scaleWithinTenMITHeavy}{42\,\%\xspace}                      % share of mix designs with the heaviest class predicted within +-10 %, MIT (2-minute instantaneous test)
\newcommand{\scaleMaeMST}{23\,\%\xspace}                                 % mean absolute error of the harmonic prediction per mix design, MST (30-minute sustainable test)
\newcommand{\scaleBiasDesignMST}{\ensuremath{-}3.1\,\%\xspace}           % mean departure from the harmonic law per mix design, MST (30-minute sustainable test)
\newcommand{\scaleWithinTenMST}{48\,\%\xspace}                           % share of mix designs predicted within +-10 %, MST (30-minute sustainable test)
\newcommand{\scaleWithinTenMSTLight}{58\,\%\xspace}                      % share of mix designs without the heaviest class predicted within +-10 %, MST (30-minute sustainable test)
\newcommand{\scaleWithinTenMSTHeavy}{25\,\%\xspace}                      % share of mix designs with the heaviest class predicted within +-10 %, MST (30-minute sustainable test)
\newcommand{\scaleMitMstMeanP}{0.507\xspace}                             % p-value (Mann-Whitney) for a difference in mean bias between MIT and MST designs
\newcommand{\scaleMitMstLeveneP}{0.0005\xspace}                          % p-value (Levene) for a difference in spread of the errors between MIT and MST
\newcommand{\scaleMitMstAbsP}{0.020\xspace}                              % p-value (Mann-Whitney) for a difference in absolute error between MIT and MST
\newcommand{\scaleMitMstLeveneLightP}{0.007\xspace}                      % Levene p-value for MIT vs MST spread, mixes without the heaviest class only
\newcommand{\scaleModelEffectP}{0.0007\xspace}                           % ANOVA p-value of the model effect on the departure from additivity (controlling test and heavy class)
\newcommand{\scaleModelEffectF}{7.8\xspace}                              % ANOVA F statistic of the model effect on the departure from additivity
\newcommand{\scaleKindEffectP}{0.783\xspace}                             % ANOVA p-value of the test (MIT/MST) effect on the departure from additivity
\newcommand{\scaleModelEffectLightP}{0.544\xspace}                       % ANOVA p-value of the model effect, mixes without the heaviest class only
\newcommand{\scaleFamilyKwP}{0.849\xspace}                               % Kruskal-Wallis p-value for differences between mix-design families (equal shares, by cost, 4-5 classes)
\newcommand{\scaleHeavyDeepSeekMIT}{\ensuremath{-}12\,\%\xspace}         % mean departure from the harmonic law of DeepSeek-R1-Distill-Qwen-7B mix designs containing the heaviest class (600-1000 in, 1000-1500 out), MIT (2-minute instantaneous test)
\newcommand{\scaleHeavyDeepSeekMST}{\ensuremath{-}8\,\%\xspace}          % mean departure from the harmonic law of DeepSeek-R1-Distill-Qwen-7B mix designs containing the heaviest class (600-1000 in, 1000-1500 out), MST (30-minute sustainable test)
\newcommand{\scaleHeavyGemmaMIT}{+16\,\%\xspace}                         % mean departure from the harmonic law of Gemma-7B mix designs containing the heaviest class (600-1000 in, 1000-1500 out), MIT (2-minute instantaneous test)
\newcommand{\scaleHeavyGemmaMST}{+64\,\%\xspace}                         % mean departure from the harmonic law of Gemma-7B mix designs containing the heaviest class (600-1000 in, 1000-1500 out), MST (30-minute sustainable test)
\newcommand{\scaleHeavyLlamaMIT}{\ensuremath{-}9\,\%\xspace}             % mean departure from the harmonic law of Llama-3.1-8B-Instruct mix designs containing the heaviest class (600-1000 in, 1000-1500 out), MIT (2-minute instantaneous test)
\newcommand{\scaleHeavyLlamaMST}{\ensuremath{-}24\,\%\xspace}            % mean departure from the harmonic law of Llama-3.1-8B-Instruct mix designs containing the heaviest class (600-1000 in, 1000-1500 out), MST (30-minute sustainable test)
\newcommand{\scaleBurstMedian}{3.50\xspace}                              % median burst factor B = MIT/MST over all request classes and models (frequentist)
\newcommand{\scaleBurstQOne}{3.41\xspace}                                % first quartile of the burst factor over all classes
\newcommand{\scaleBurstQThree}{3.94\xspace}                              % third quartile of the burst factor over all classes
\newcommand{\scaleBurstOnBand}{77\,\%\xspace}                            % share of request classes whose B is within +-15 % of their model's median
\newcommand{\scaleBurstNCells}{60\xspace}                                % number of request classes (cells) for which B was measured (3 models x 20 classes)
\newcommand{\scaleBurstDeepSeek}{3.57\xspace}                            % median burst factor B = MIT/MST of DeepSeek-R1-Distill-Qwen-7B over its 20 classes (Bayesian when available)
\newcommand{\scaleBurstMaxDeepSeek}{5.6\xspace}                          % largest burst factor among DeepSeek-R1-Distill-Qwen-7B's classes
\newcommand{\scaleBurstOnBandDeepSeek}{70\,\%\xspace}                    % share of DeepSeek-R1-Distill-Qwen-7B's classes with B within +-15 % of its median
\newcommand{\scaleBurstDeepSeekLo}{3.49\xspace}                          % lower bound of the 95 % credible interval of DeepSeek-R1-Distill-Qwen-7B's median burst factor
\newcommand{\scaleBurstDeepSeekHi}{3.64\xspace}                          % upper bound of the 95 % credible interval of DeepSeek-R1-Distill-Qwen-7B's median burst factor
\newcommand{\scaleBurstGemma}{3.56\xspace}                               % median burst factor B = MIT/MST of Gemma-7B over its 20 classes (Bayesian when available)
\newcommand{\scaleBurstMaxGemma}{11.6\xspace}                            % largest burst factor among Gemma-7B's classes
\newcommand{\scaleBurstOnBandGemma}{75\,\%\xspace}                       % share of Gemma-7B's classes with B within +-15 % of its median
\newcommand{\scaleBurstGemmaLo}{3.49\xspace}                             % lower bound of the 95 % credible interval of Gemma-7B's median burst factor
\newcommand{\scaleBurstGemmaHi}{3.65\xspace}                             % upper bound of the 95 % credible interval of Gemma-7B's median burst factor
\newcommand{\scaleBurstLlama}{3.55\xspace}                               % median burst factor B = MIT/MST of Llama-3.1-8B-Instruct over its 20 classes (Bayesian when available)
\newcommand{\scaleBurstMaxLlama}{5.4\xspace}                             % largest burst factor among Llama-3.1-8B-Instruct's classes
\newcommand{\scaleBurstOnBandLlama}{85\,\%\xspace}                       % share of Llama-3.1-8B-Instruct's classes with B within +-15 % of its median
\newcommand{\scaleBurstLlamaLo}{3.48\xspace}                             % lower bound of the 95 % credible interval of Llama-3.1-8B-Instruct's median burst factor
\newcommand{\scaleBurstLlamaHi}{3.62\xspace}                             % upper bound of the 95 % credible interval of Llama-3.1-8B-Instruct's median burst factor
\newcommand{\scaleBurstModelP}{0.742\xspace}                             % ANOVA p-value of the model effect on log B (controlling input and output range)
\newcommand{\scaleBurstOutP}{0.139\xspace}                               % ANOVA p-value of the output-length effect on log B
\newcommand{\scaleBurstInP}{0.672\xspace}                                % ANOVA p-value of the input-length effect on log B
\newcommand{\scaleBurstOutRobustP}{\ensuremath{6\times 10^{-7}}\xspace}  % ANOVA p-value of the output-length effect on log B without Gemma's two anomalous cells
\newcommand{\scaleBurstModelRobustP}{0.045\xspace}                       % ANOVA p-value of the model effect on log B without Gemma's two anomalous cells
\newcommand{\scaleBurstShort}{4.15\xspace}                               % geometric-mean burst factor of classes with short answers (1-100 output tokens)
\newcommand{\scaleBurstRest}{3.68\xspace}                                % geometric-mean burst factor of the other classes
\newcommand{\scaleBurstShortRatio}{+13\,\%\xspace}                       % how much larger B is for short answers than for the rest
\newcommand{\scaleBurstShortP}{0.0004\xspace}                            % Mann-Whitney p-value for short answers vs the rest
\newcommand{\scaleBurstPaperExample}{1.4\xspace}                         % burst factor implied by the paper's introductory example (1,400 req/min vs 30,000 req/30 min)
\newcommand{\scaleBurstMixN}{12\xspace}                                  % number of mix designs measured with the same composition in both MIT and MST
\newcommand{\scaleBurstMixErr}{+19\,\%\xspace}                           % mean excess of the measured burst factor of mixes over the one predicted from their classes
\newcommand{\scaleBurstMixWithin}{50\,\%\xspace}                         % share of those mixes whose measured B is within +-15 % of the prediction
\newcommand{\scaleBurstEquivMin}{100\,\%\xspace}                         % smallest posterior probability, over model pairs, that two models' burst factors are equivalent within +-10 %
\newcommand{\scaleBurstPairMaxDiff}{0.5\,\%\xspace}                      % largest posterior-median difference between the burst factors of two models
\newcommand{\scaleBurstPairCi}{3.5\,\%\xspace}                           % largest distance from 1 of the 95 % credible bounds of the ratio between two models' burst factors
\newcommand{\scaleTbWindow}{2\xspace}                                    % MIT window W in minutes used to characterise bursts
\newcommand{\scaleTbDepth}{5.0\xspace}                                   % token-bucket depth b/s in minutes of refill, b = (B-1) W, with the typical B
\newcommand{\scaleTbDepthSafe}{4.6\xspace}                               % token-bucket depth in minutes of refill with the conservative B_safe
\newcommand{\scaleTbDepthSmall}{6.1\xspace}                              % token-bucket depth in minutes of refill for plans of requests up to 300 tokens
\newcommand{\scaleTbBSafe}{3.29\xspace}                                  % conservative burst factor B_safe: 10th percentile over classes (covers 90 % of them)
\newcommand{\scaleTbBSmall}{4.05\xspace}                                 % median burst factor of classes up to 300 input and output tokens (the paper's plans)
\newcommand{\scaleTbRho}{1.2\xspace}                                     % oversubscription ratio rho used in the paper's example
\newcommand{\scaleTbPhiMax}{77\,\%\xspace}                               % largest share of the allocation that may burst at once at rho = TbRho: phi <= (B/rho - 1)/(B - 1)
\newcommand{\scaleTbRhoMaxHalf}{1.56\xspace}                             % largest burst-safe oversubscription ratio if half of the allocation bursts at once: rho <= B/(1 + phi(B-1))
\newcommand{\scaleTbRhoMaxQuarter}{2.15\xspace}                          % largest burst-safe oversubscription ratio if a quarter of the allocation bursts at once
\newcommand{\scaleTbPaperRatio}{21\xspace}                               % burst-to-sustained rate ratio implied by Table 3 of the paper (Basic plan)
\newcommand{\scaleTbSkMaxDeepSeek}{2{,}307\xspace}                       % SK_max of DeepSeek-R1-Distill-Qwen-7B: largest MST in its capacity matrix, req/min (1 K per request of that class)
\newcommand{\scaleTbThrottleDeepSeek}{8{,}226\xspace}                    % throttling threshold Theta = B * SK_max of DeepSeek-R1-Distill-Qwen-7B, sustainable units per minute averaged over W
\newcommand{\scaleTbSkMaxGemma}{2{,}047\xspace}                          % SK_max of Gemma-7B: largest MST in its capacity matrix, req/min (1 K per request of that class)
\newcommand{\scaleTbThrottleGemma}{7{,}297\xspace}                       % throttling threshold Theta = B * SK_max of Gemma-7B, sustainable units per minute averaged over W
\newcommand{\scaleTbSkMaxLlama}{2{,}176\xspace}                          % SK_max of Llama-3.1-8B-Instruct: largest MST in its capacity matrix, req/min (1 K per request of that class)
\newcommand{\scaleTbThrottleLlama}{7{,}717\xspace}                       % throttling threshold Theta = B * SK_max of Llama-3.1-8B-Instruct, sustainable units per minute averaged over W
\newcommand{\scaleTbBasicBurstTwo}{14\xspace}                            % burst allowance of the paper's Basic plan (50 req/30 min) derived from B, in requests per 2 minutes
\newcommand{\scaleTbBasicDepth}{10\xspace}                               % token-bucket depth of the Basic plan derived from B, in requests
\newcommand{\scaleTbProBurstTwo}{47\xspace}                              % burst allowance of the paper's Pro plan (175 req/30 min) derived from B, in requests per 2 minutes
\newcommand{\scaleTbProDepth}{36\xspace}                                 % token-bucket depth of the Pro plan derived from B, in requests
 in the preamble; every macro is prefixed with \scale.
% Each line ends with a comment saying what the value means, measures or represents
% (descriptions live in analysis/macro_docs.py). Percentages are relative differences
% unless stated; 'design level' means one value per mix design, repeats averaged.
\providecommand{\xspace}{}
\newcommand{\scaleNEvaluations}{7{,}944\xspace}                          % number of MoST pass/fail SLA evaluations analysed (all runs, after removing exact copies)
\newcommand{\scaleNMixRuns}{301\xspace}                                  % number of mix runs (MoST searches on request mixes)
\newcommand{\scaleNDesigns}{114\xspace}                                  % number of distinct mix designs (model x test x composition); unit of the frequentist tests
\newcommand{\scaleNRepeatRuns}{270\xspace}                               % number of mix runs that belong to designs measured more than once
\newcommand{\scaleNRepeatDesigns}{83\xspace}                             % number of mix designs measured more than once (stability check)
\newcommand{\scaleNDuplicates}{114\xspace}                               % number of byte-identical copies of a run found in the archive and counted once
\newcommand{\scaleNPureRuns}{360\xspace}                                 % number of pure-class runs (one request class per MoST search)
\newcommand{\scaleQHatFreq}{\ensuremath{-}0.91\xspace}                   % power-mean exponent q of the combination law, profile-likelihood estimate (q=-1 harmonic, 0 geometric, +1 arithmetic)
\newcommand{\scaleQFreqLo}{\ensuremath{-}0.98\xspace}                    % lower bound of the 95 % profile-likelihood confidence interval of q
\newcommand{\scaleQFreqHi}{\ensuremath{-}0.84\xspace}                    % upper bound of the 95 % profile-likelihood confidence interval of q
\newcommand{\scaleHarmonicWinShare}{97\,\%\xspace}                       % share of mix runs for which the harmonic prediction is closer to the measurement than the arithmetic one
\newcommand{\scaleCalibSlope}{0.96\xspace}                               % slope of the regression log(measured mix capacity) ~ log(harmonic prediction); 1 = correct scaling
\newcommand{\scaleCalibSlopeLo}{0.90\xspace}                             % lower bound of the 95 % confidence interval of the calibration slope (cluster-robust by model)
\newcommand{\scaleCalibSlopeHi}{1.02\xspace}                             % upper bound of the 95 % confidence interval of the calibration slope
\newcommand{\scaleCalibRsq}{0.97\xspace}                                 % R^2 of the calibration regression on the log scale
\newcommand{\scaleBiasFreq}{\ensuremath{-}3.0\,\%\xspace}                % mean departure of measured mix capacity from the harmonic prediction, frequentist, designs as units
\newcommand{\scaleBiasFreqLo}{\ensuremath{-}6.9\,\%\xspace}              % lower bound of the 95 % confidence interval of BiasFreq
\newcommand{\scaleBiasFreqHi}{+1.0\,\%\xspace}                           % upper bound of the 95 % confidence interval of BiasFreq
\newcommand{\scaleBiasTostP}{0.001\xspace}                               % p-value of the TOST equivalence test of the bias against +-10 % (small = practically zero)
\newcommand{\scaleArithMeanErr}{\ensuremath{-}53\,\%\xspace}             % mean error of the arithmetic-mean law (measured vs predicted mix capacity, design level); negative = overestimates
\newcommand{\scaleGeomMeanErr}{\ensuremath{-}39\,\%\xspace}              % mean error of the geometric-mean law (design level); negative = overestimates
\newcommand{\scaleBoundsIn}{294\xspace}                                  % number of mix runs whose capacity lies between their slowest and fastest class
\newcommand{\scaleBoundsN}{301\xspace}                                   % number of mix runs checked for boundedness
\newcommand{\scaleBoundsShare}{97.7\,\%\xspace}                          % share of mix runs within the capacity range of their classes
\newcommand{\scaleAboveArith}{0\xspace}                                  % number of mix runs whose capacity exceeds the arithmetic mean of their classes
\newcommand{\scaleQHatBayes}{\ensuremath{-}0.96\xspace}                  % power-mean exponent q, Bayesian posterior median (free-q model)
\newcommand{\scaleQBayesLo}{\ensuremath{-}1.02\xspace}                   % lower bound of the 95 % credible interval of q
\newcommand{\scaleQBayesHi}{\ensuremath{-}0.88\xspace}                   % upper bound of the 95 % credible interval of q
\newcommand{\scaleBiasBayes}{\ensuremath{-}5.5\,\%\xspace}               % pooled bias delta_0 of mixes w.r.t. the harmonic law, Bayesian posterior median
\newcommand{\scaleBiasBayesLo}{\ensuremath{-}9.9\,\%\xspace}             % lower bound of the 90 % credible interval of delta_0
\newcommand{\scaleBiasBayesHi}{\ensuremath{-}1.0\,\%\xspace}             % upper bound of the 90 % credible interval of delta_0
\newcommand{\scaleBiasBayesPRope}{92\,\%\xspace}                         % posterior probability that |delta_0| < 10 % (region of practical equivalence)
\newcommand{\scalePredFree}{+1.7\xspace}                                 % out-of-sample log predictive density of the free-q law minus that of the harmonic law (higher = better)
\newcommand{\scalePredFreeSE}{3.3\xspace}                                % standard error of PredFree
\newcommand{\scalePredGeom}{\ensuremath{-}264.5\xspace}                  % out-of-sample log predictive density of the geometric law minus that of the harmonic law
\newcommand{\scalePredGeomSE}{24.8\xspace}                               % standard error of PredGeom
\newcommand{\scalePredArith}{\ensuremath{-}441.4\xspace}                 % out-of-sample log predictive density of the arithmetic law minus that of the harmonic law
\newcommand{\scalePredArithSE}{38.1\xspace}                              % standard error of PredArith
\newcommand{\scaleSdMixMIT}{6\,\%\xspace}                                % posterior SD of the design effect (departure of a mix design from the harmonic law), MIT (2-minute instantaneous test), as %
\newcommand{\scaleSdRunMIT}{0.3\,\%\xspace}                              % posterior SD between repeated runs of the same mix design, MIT (2-minute instantaneous test), as %
\newcommand{\scaleSdRepMIT}{0.3\,\%\xspace}                              % posterior SD between repeated runs of the same pure class, MIT (2-minute instantaneous test), as %
\newcommand{\scaleBandMIT}{\ensuremath{\pm}15\,\%\xspace}                % half-width of the 90 % planning band around the harmonic prediction, MIT (2-minute instantaneous test)
\newcommand{\scaleSdMixMST}{21\,\%\xspace}                               % posterior SD of the design effect (departure of a mix design from the harmonic law), MST (30-minute sustainable test), as %
\newcommand{\scaleSdRunMST}{1.8\,\%\xspace}                              % posterior SD between repeated runs of the same mix design, MST (30-minute sustainable test), as %
\newcommand{\scaleSdRepMST}{1.7\,\%\xspace}                              % posterior SD between repeated runs of the same pure class, MST (30-minute sustainable test), as %
\newcommand{\scaleBandMST}{\ensuremath{\pm}39\,\%\xspace}                % half-width of the 90 % planning band around the harmonic prediction, MST (30-minute sustainable test)
\newcommand{\scaleBiasDeepSeekMIT}{\ensuremath{-}8.8\,\%\xspace}         % Bayesian bias delta of DeepSeek-R1-Distill-Qwen-7B mixes w.r.t. the harmonic law, MIT (2-minute instantaneous test)
\newcommand{\scaleBiasDeepSeekMST}{\ensuremath{-}4.1\,\%\xspace}         % Bayesian bias delta of DeepSeek-R1-Distill-Qwen-7B mixes w.r.t. the harmonic law, MST (30-minute sustainable test)
\newcommand{\scaleBiasGemmaMIT}{\ensuremath{-}0.3\,\%\xspace}            % Bayesian bias delta of Gemma-7B mixes w.r.t. the harmonic law, MIT (2-minute instantaneous test)
\newcommand{\scaleBiasGemmaMST}{\ensuremath{-}4.6\,\%\xspace}            % Bayesian bias delta of Gemma-7B mixes w.r.t. the harmonic law, MST (30-minute sustainable test)
\newcommand{\scaleBiasLlamaMIT}{\ensuremath{-}7.0\,\%\xspace}            % Bayesian bias delta of Llama-3.1-8B-Instruct mixes w.r.t. the harmonic law, MIT (2-minute instantaneous test)
\newcommand{\scaleBiasLlamaMST}{\ensuremath{-}7.8\,\%\xspace}            % Bayesian bias delta of Llama-3.1-8B-Instruct mixes w.r.t. the harmonic law, MST (30-minute sustainable test)
\newcommand{\scaleBayesDraws}{600\xspace}                                % NUTS draws per chain (4 chains, same number of tuning steps)
\newcommand{\scaleMaxRhat}{1.02\xspace}                                  % largest split-R-hat over the global parameters of the reported Bayesian models (<1.05 = converged)
\newcommand{\scaleMaeMIT}{10\,\%\xspace}                                 % mean absolute error of the harmonic prediction per mix design, MIT (2-minute instantaneous test)
\newcommand{\scaleBiasDesignMIT}{\ensuremath{-}2.9\,\%\xspace}           % mean departure from the harmonic law per mix design, MIT (2-minute instantaneous test)
\newcommand{\scaleWithinTenMIT}{61\,\%\xspace}                           % share of mix designs predicted within +-10 %, MIT (2-minute instantaneous test)
\newcommand{\scaleWithinTenMITLight}{75\,\%\xspace}                      % share of mix designs without the heaviest class predicted within +-10 %, MIT (2-minute instantaneous test)
\newcommand{\scaleWithinTenMITHeavy}{42\,\%\xspace}                      % share of mix designs with the heaviest class predicted within +-10 %, MIT (2-minute instantaneous test)
\newcommand{\scaleMaeMST}{23\,\%\xspace}                                 % mean absolute error of the harmonic prediction per mix design, MST (30-minute sustainable test)
\newcommand{\scaleBiasDesignMST}{\ensuremath{-}3.1\,\%\xspace}           % mean departure from the harmonic law per mix design, MST (30-minute sustainable test)
\newcommand{\scaleWithinTenMST}{48\,\%\xspace}                           % share of mix designs predicted within +-10 %, MST (30-minute sustainable test)
\newcommand{\scaleWithinTenMSTLight}{58\,\%\xspace}                      % share of mix designs without the heaviest class predicted within +-10 %, MST (30-minute sustainable test)
\newcommand{\scaleWithinTenMSTHeavy}{25\,\%\xspace}                      % share of mix designs with the heaviest class predicted within +-10 %, MST (30-minute sustainable test)
\newcommand{\scaleMitMstMeanP}{0.507\xspace}                             % p-value (Mann-Whitney) for a difference in mean bias between MIT and MST designs
\newcommand{\scaleMitMstLeveneP}{0.0005\xspace}                          % p-value (Levene) for a difference in spread of the errors between MIT and MST
\newcommand{\scaleMitMstAbsP}{0.020\xspace}                              % p-value (Mann-Whitney) for a difference in absolute error between MIT and MST
\newcommand{\scaleMitMstLeveneLightP}{0.007\xspace}                      % Levene p-value for MIT vs MST spread, mixes without the heaviest class only
\newcommand{\scaleModelEffectP}{0.0007\xspace}                           % ANOVA p-value of the model effect on the departure from additivity (controlling test and heavy class)
\newcommand{\scaleModelEffectF}{7.8\xspace}                              % ANOVA F statistic of the model effect on the departure from additivity
\newcommand{\scaleKindEffectP}{0.783\xspace}                             % ANOVA p-value of the test (MIT/MST) effect on the departure from additivity
\newcommand{\scaleModelEffectLightP}{0.544\xspace}                       % ANOVA p-value of the model effect, mixes without the heaviest class only
\newcommand{\scaleFamilyKwP}{0.849\xspace}                               % Kruskal-Wallis p-value for differences between mix-design families (equal shares, by cost, 4-5 classes)
\newcommand{\scaleHeavyDeepSeekMIT}{\ensuremath{-}12\,\%\xspace}         % mean departure from the harmonic law of DeepSeek-R1-Distill-Qwen-7B mix designs containing the heaviest class (600-1000 in, 1000-1500 out), MIT (2-minute instantaneous test)
\newcommand{\scaleHeavyDeepSeekMST}{\ensuremath{-}8\,\%\xspace}          % mean departure from the harmonic law of DeepSeek-R1-Distill-Qwen-7B mix designs containing the heaviest class (600-1000 in, 1000-1500 out), MST (30-minute sustainable test)
\newcommand{\scaleHeavyGemmaMIT}{+16\,\%\xspace}                         % mean departure from the harmonic law of Gemma-7B mix designs containing the heaviest class (600-1000 in, 1000-1500 out), MIT (2-minute instantaneous test)
\newcommand{\scaleHeavyGemmaMST}{+64\,\%\xspace}                         % mean departure from the harmonic law of Gemma-7B mix designs containing the heaviest class (600-1000 in, 1000-1500 out), MST (30-minute sustainable test)
\newcommand{\scaleHeavyLlamaMIT}{\ensuremath{-}9\,\%\xspace}             % mean departure from the harmonic law of Llama-3.1-8B-Instruct mix designs containing the heaviest class (600-1000 in, 1000-1500 out), MIT (2-minute instantaneous test)
\newcommand{\scaleHeavyLlamaMST}{\ensuremath{-}24\,\%\xspace}            % mean departure from the harmonic law of Llama-3.1-8B-Instruct mix designs containing the heaviest class (600-1000 in, 1000-1500 out), MST (30-minute sustainable test)
\newcommand{\scaleBurstMedian}{3.50\xspace}                              % median burst factor B = MIT/MST over all request classes and models (frequentist)
\newcommand{\scaleBurstQOne}{3.41\xspace}                                % first quartile of the burst factor over all classes
\newcommand{\scaleBurstQThree}{3.94\xspace}                              % third quartile of the burst factor over all classes
\newcommand{\scaleBurstOnBand}{77\,\%\xspace}                            % share of request classes whose B is within +-15 % of their model's median
\newcommand{\scaleBurstNCells}{60\xspace}                                % number of request classes (cells) for which B was measured (3 models x 20 classes)
\newcommand{\scaleBurstDeepSeek}{3.57\xspace}                            % median burst factor B = MIT/MST of DeepSeek-R1-Distill-Qwen-7B over its 20 classes (Bayesian when available)
\newcommand{\scaleBurstMaxDeepSeek}{5.6\xspace}                          % largest burst factor among DeepSeek-R1-Distill-Qwen-7B's classes
\newcommand{\scaleBurstOnBandDeepSeek}{70\,\%\xspace}                    % share of DeepSeek-R1-Distill-Qwen-7B's classes with B within +-15 % of its median
\newcommand{\scaleBurstDeepSeekLo}{3.49\xspace}                          % lower bound of the 95 % credible interval of DeepSeek-R1-Distill-Qwen-7B's median burst factor
\newcommand{\scaleBurstDeepSeekHi}{3.64\xspace}                          % upper bound of the 95 % credible interval of DeepSeek-R1-Distill-Qwen-7B's median burst factor
\newcommand{\scaleBurstGemma}{3.56\xspace}                               % median burst factor B = MIT/MST of Gemma-7B over its 20 classes (Bayesian when available)
\newcommand{\scaleBurstMaxGemma}{11.6\xspace}                            % largest burst factor among Gemma-7B's classes
\newcommand{\scaleBurstOnBandGemma}{75\,\%\xspace}                       % share of Gemma-7B's classes with B within +-15 % of its median
\newcommand{\scaleBurstGemmaLo}{3.49\xspace}                             % lower bound of the 95 % credible interval of Gemma-7B's median burst factor
\newcommand{\scaleBurstGemmaHi}{3.65\xspace}                             % upper bound of the 95 % credible interval of Gemma-7B's median burst factor
\newcommand{\scaleBurstLlama}{3.55\xspace}                               % median burst factor B = MIT/MST of Llama-3.1-8B-Instruct over its 20 classes (Bayesian when available)
\newcommand{\scaleBurstMaxLlama}{5.4\xspace}                             % largest burst factor among Llama-3.1-8B-Instruct's classes
\newcommand{\scaleBurstOnBandLlama}{85\,\%\xspace}                       % share of Llama-3.1-8B-Instruct's classes with B within +-15 % of its median
\newcommand{\scaleBurstLlamaLo}{3.48\xspace}                             % lower bound of the 95 % credible interval of Llama-3.1-8B-Instruct's median burst factor
\newcommand{\scaleBurstLlamaHi}{3.62\xspace}                             % upper bound of the 95 % credible interval of Llama-3.1-8B-Instruct's median burst factor
\newcommand{\scaleBurstModelP}{0.742\xspace}                             % ANOVA p-value of the model effect on log B (controlling input and output range)
\newcommand{\scaleBurstOutP}{0.139\xspace}                               % ANOVA p-value of the output-length effect on log B
\newcommand{\scaleBurstInP}{0.672\xspace}                                % ANOVA p-value of the input-length effect on log B
\newcommand{\scaleBurstOutRobustP}{\ensuremath{6\times 10^{-7}}\xspace}  % ANOVA p-value of the output-length effect on log B without Gemma's two anomalous cells
\newcommand{\scaleBurstModelRobustP}{0.045\xspace}                       % ANOVA p-value of the model effect on log B without Gemma's two anomalous cells
\newcommand{\scaleBurstShort}{4.15\xspace}                               % geometric-mean burst factor of classes with short answers (1-100 output tokens)
\newcommand{\scaleBurstRest}{3.68\xspace}                                % geometric-mean burst factor of the other classes
\newcommand{\scaleBurstShortRatio}{+13\,\%\xspace}                       % how much larger B is for short answers than for the rest
\newcommand{\scaleBurstShortP}{0.0004\xspace}                            % Mann-Whitney p-value for short answers vs the rest
\newcommand{\scaleBurstPaperExample}{1.4\xspace}                         % burst factor implied by the paper's introductory example (1,400 req/min vs 30,000 req/30 min)
\newcommand{\scaleBurstMixN}{12\xspace}                                  % number of mix designs measured with the same composition in both MIT and MST
\newcommand{\scaleBurstMixErr}{+19\,\%\xspace}                           % mean excess of the measured burst factor of mixes over the one predicted from their classes
\newcommand{\scaleBurstMixWithin}{50\,\%\xspace}                         % share of those mixes whose measured B is within +-15 % of the prediction
\newcommand{\scaleBurstEquivMin}{100\,\%\xspace}                         % smallest posterior probability, over model pairs, that two models' burst factors are equivalent within +-10 %
\newcommand{\scaleBurstPairMaxDiff}{0.5\,\%\xspace}                      % largest posterior-median difference between the burst factors of two models
\newcommand{\scaleBurstPairCi}{3.5\,\%\xspace}                           % largest distance from 1 of the 95 % credible bounds of the ratio between two models' burst factors
\newcommand{\scaleTbWindow}{2\xspace}                                    % MIT window W in minutes used to characterise bursts
\newcommand{\scaleTbDepth}{5.0\xspace}                                   % token-bucket depth b/s in minutes of refill, b = (B-1) W, with the typical B
\newcommand{\scaleTbDepthSafe}{4.6\xspace}                               % token-bucket depth in minutes of refill with the conservative B_safe
\newcommand{\scaleTbDepthSmall}{6.1\xspace}                              % token-bucket depth in minutes of refill for plans of requests up to 300 tokens
\newcommand{\scaleTbBSafe}{3.29\xspace}                                  % conservative burst factor B_safe: 10th percentile over classes (covers 90 % of them)
\newcommand{\scaleTbBSmall}{4.05\xspace}                                 % median burst factor of classes up to 300 input and output tokens (the paper's plans)
\newcommand{\scaleTbRho}{1.2\xspace}                                     % oversubscription ratio rho used in the paper's example
\newcommand{\scaleTbPhiMax}{77\,\%\xspace}                               % largest share of the allocation that may burst at once at rho = TbRho: phi <= (B/rho - 1)/(B - 1)
\newcommand{\scaleTbRhoMaxHalf}{1.56\xspace}                             % largest burst-safe oversubscription ratio if half of the allocation bursts at once: rho <= B/(1 + phi(B-1))
\newcommand{\scaleTbRhoMaxQuarter}{2.15\xspace}                          % largest burst-safe oversubscription ratio if a quarter of the allocation bursts at once
\newcommand{\scaleTbPaperRatio}{21\xspace}                               % burst-to-sustained rate ratio implied by Table 3 of the paper (Basic plan)
\newcommand{\scaleTbSkMaxDeepSeek}{2{,}307\xspace}                       % SK_max of DeepSeek-R1-Distill-Qwen-7B: largest MST in its capacity matrix, req/min (1 K per request of that class)
\newcommand{\scaleTbThrottleDeepSeek}{8{,}226\xspace}                    % throttling threshold Theta = B * SK_max of DeepSeek-R1-Distill-Qwen-7B, sustainable units per minute averaged over W
\newcommand{\scaleTbSkMaxGemma}{2{,}047\xspace}                          % SK_max of Gemma-7B: largest MST in its capacity matrix, req/min (1 K per request of that class)
\newcommand{\scaleTbThrottleGemma}{7{,}297\xspace}                       % throttling threshold Theta = B * SK_max of Gemma-7B, sustainable units per minute averaged over W
\newcommand{\scaleTbSkMaxLlama}{2{,}176\xspace}                          % SK_max of Llama-3.1-8B-Instruct: largest MST in its capacity matrix, req/min (1 K per request of that class)
\newcommand{\scaleTbThrottleLlama}{7{,}717\xspace}                       % throttling threshold Theta = B * SK_max of Llama-3.1-8B-Instruct, sustainable units per minute averaged over W
\newcommand{\scaleTbBasicBurstTwo}{14\xspace}                            % burst allowance of the paper's Basic plan (50 req/30 min) derived from B, in requests per 2 minutes
\newcommand{\scaleTbBasicDepth}{10\xspace}                               % token-bucket depth of the Basic plan derived from B, in requests
\newcommand{\scaleTbProBurstTwo}{47\xspace}                              % burst allowance of the paper's Pro plan (175 req/30 min) derived from B, in requests per 2 minutes
\newcommand{\scaleTbProDepth}{36\xspace}                                 % token-bucket depth of the Pro plan derived from B, in requests
